\chapter{Computational Acceleration for Large-Q Regime}
\label{cha:computational-acceleration}

Chapter~\ref{cha:methodology} derived a correct and convergent Laplace-EM algorithm, but left its dominant per-iteration cost unaddressed. Each Newton step in E-step~\ref{sec: E-step} requires forming and solving the $Q\times Q$ system~\ref{eq: hessian-equality} at cost $\mathcal{O}(N_{j}Q^{2}+Q^{3})$, in the meanwhile, each inner iteration of the M-step factor update~\ref{eq: Rubin3}-\ref{eq: Rubin4} requires forming $\bm{\Sigma}^{-1}$ densely at cost $\mathcal{O}(Q^{3})$. For large Q, these costs are not merely inconvenient but make the algorithm impractical. This chapter develops an exact acceleration method to remove this bottleneck by showing that the Poisson-Multinomial equivalence that already underlies the M-step can be redeployed via auxiliary-variable construction, which can remove the cross-category coupling from the E-step Hessian as well. 

\section{PME Auxiliary-Variable Construction}
\label{sec:auxiliary-variable-construction}

For a fixed group $j$, we introduce an auxiliary scalar $\xi_{i}$ per observation $i\in \mathcal{I}_{j}$ and define the Poisson surrogate log-posterior
\begin{equation}
    \ell_P(\boldsymbol{\lambda}, \boldsymbol{\xi}) \;=\; \sum_{i \in \mathcal{I}_j} \sum_{q=1}^{Q} \Big[ y_{iq}\, \eta_{iq}(\boldsymbol{\lambda}) - \exp\big(\eta_{iq}(\boldsymbol{\lambda}) + \xi_i\big) \Big] \;-\; \tfrac{1}{2} \boldsymbol{\lambda}^\top \boldsymbol{\Sigma}^{-1} \boldsymbol{\lambda},
\end{equation}
which treats the counts as independent Poisson draws with rate $\exp{(\eta_{iq}+\xi_{i})}$ rather than as a multinomial draw. This is the same auxiliary-variable idea that underlies Proposition~\ref{prop:PME}, but it is now applied to the E-step posterior instead of M-step likelihood.

\section{Diagonal Low Rank Hessian \& Woodbury Solve}
\label{sec:diagonal-low-rank-hessian}

With $\bm{\xi}$ fixed at $\bm{\xi}^{*}(\bm{\lambda})$, we denote $d_{q}:=\sum_{i\in \mathcal{I}_{j}}\exp{\bigl(\eta_{iq}(\bm{\lambda})+\xi_{i}\bigr)}$. The Hessian of $\ell_{P}(\cdot, \bm{\xi})$ with respect to $\bm{\lambda}$ is 
\begin{equation}
    \mathbf{H}_{P}(\bm{\lambda})=-\text{diag}(\mathbf{d})-\bm{\Sigma}^{-1}, \quad \mathbf{d}=(d_{1},\dots, d_{Q})^{\top},
    \label{eq:acceleration-hessian1}
\end{equation}
which is purely diagonal in its Poisson part. We expand $\bm{\Sigma}^{-1}$ by the Woodbury identity

\begin{equation}
    \bm{\Sigma}^{-1}=\sigma^{-2}\mathbf{I}_{Q}-\sigma^{-4}\mathbf{BM}_{K}^{-1}\mathbf{B}^{\top} \text{ with } \mathbf{M}_{K}=\mathbf{I}_{K}+\sigma^{-2}\mathbf{B}^{\top}\mathbf{B}.
    \label{eq:woodbury2}
\end{equation}

We write $\tilde{\mathbf{D}}=\text{diag}(\mathbf{d}+\sigma^{-2}\mathbf{1}_{Q})$ for the combined diagonal term
\begin{equation}
    -\mathbf{H}_{P}(\bm{\lambda})=\tilde{\mathbf{D}}-\sigma^{-4}\mathbf{BM}_{K}^{-1}\mathbf{B}^{\top},
    \label{eq:acceleration-hessian3}
\end{equation}
it's a $Q\times Q$ matrix. Applying Woodbury identity once more and we obtain
\begin{equation}
    \bigl(-\mathbf{H}_{P}(\bm{\lambda})\bigr)^{-1}=\tilde{\mathbf{D}}^{-1}-\tilde{\mathbf{D}}^{-1}\mathbf{B}\bm{\Gamma}^{-1}\mathbf{B}^{\top}\tilde{\mathbf{D}}^{-1} \text{ with } \bm{\Gamma}:=\mathbf{B}^{\top}\tilde{\mathbf{D}}^{-1}\mathbf{B}-\sigma^{4}\mathbf{M}_{K}.
    \label{eq:inverse-hessian4}
\end{equation}

Therefore, given the Cholesky factor $\mathbf{L}$ of $-\bm{\Gamma}$, a Newton direction $\Delta\bm{\lambda}=\bigl(-\mathbf{H}_{P}(\bm{\lambda})\bigr)^{-1}\mathbf{g}$ for any gradient vector $\mathbf{g}$ is obtained as

\begin{equation}
    \Delta\boldsymbol{\lambda} \;=\; \tilde{\mathbf{D}}^{-1}\mathbf{g} \;+\; \tilde{\mathbf{D}}^{-1}\mathbf{B}(\mathbf{L}\mathbf{L}^\top)^{-1}\mathbf{B}^\top\tilde{\mathbf{D}}^{-1}\mathbf{g},
    \label{eq:newton-direction5}
\end{equation}
computed by two matrix-vector products against $\mathbf{B}$ and one $K\times K$ triangular solve, after a one-time $\mathcal{O}(QK^{2}+K^{3})$ cost to form $\bm{\Gamma}$ and its Cholesky factor.

\begin{proposition}[Negative definiteness and conditioning of $\bm{\Gamma}$]
    \label{prop:negative-definiteness}
    If $\bm{\Gamma}\prec 0$ and $\lambda_{\text{max}}(\bm{\Gamma})\leq -\sigma^{4}$, then $-\bm{\Gamma}\succ \sigma^{4}\mathbf{I}_{K}$ uniformly in $\bm{\lambda}$, $\mathbf{B}$ and data.
    \label{prop:negative-definiteness}
\end{proposition}

\begin{proof}
    Since $d_{q}\geq 0$ for each $q$, $\tilde{\mathbf{D}}=\text{diag}(\mathbf{d}+\sigma^{-2}\mathbf{1}_{Q})\succeq \sigma^{-2}\mathbf{I}_{Q}$. Therefore, $\tilde{\mathbf{D}}^{-1}\preceq \sigma^{2}\mathbf{I}_{Q}$ and $\mathbf{B}^{\top}\tilde{\mathbf{D}}^{-1}\mathbf{B}\preceq \sigma^{2}\mathbf{B}^{\top}\mathbf{B}$.
    Rewrite the definition of $\bm{\Gamma}$ in~\ref{eq:inverse-hessian4} using the substitution $\mathbf{M}_{K}=\mathbf{I}_{K}+\sigma^{-2}\mathbf{B}^{\top}\mathbf{B}$, then
    \[
        -\bm{\Gamma}=\sigma^{4}\mathbf{M}_{K}-\mathbf{B}^{\top}\tilde{\mathbf{D}}^{-1}\mathbf{B}\succeq \sigma^{4}\mathbf{M}_{K}+\sigma^{2}\mathbf{B}^{\top}\mathbf{B}-\sigma^{2}\mathbf{B}^{\top}\mathbf{B}=\sigma^{4}\mathbf{I}_{K}.
    \]
\end{proof}

Proposition~\ref{prop:negative-definiteness} guarantees that the Cholesky factorization used in~\ref{eq:newton-direction5} is always numerically well behaved. Regardless of how large $Q$ and $\mathbf{B}$ increase, the relevant $K\times K$ matrix never becomes ill-conditioned purely as an artefact of dimension, since its smallest eigenvalue is bounded below by the dimension-free quantity $\sigma^{4}$.

\section{Ascent Direction on True Posterior}
\label{sec:ascent-direction-on-true-posterior}

The direction $\Delta \bm{\lambda}$ of~\ref{eq:newton-direction5} is a Newton step for the surrogate objective $\ell_{P}(\cdot, \bm{\xi}^{*}\bigl(\bm{\lambda})\bigr)$ instead of true log-posterior $\ell$.

\begin{algorithm}[t]
    \caption{E-step inner iteration: posterior-mode search for the group effect $\bm{\lambda}_j$}
    \label{alg:estep-inner}
    \begin{algorithmic}[1]
    \Require group data $\{\bm{y}_{ij}\}_{i\in\mathcal{I}_j}$; current parameters
         $(\bm{\mu},\bm{\Phi},\mathbf{B},\sigma^2)$; Armijo constants
         $c_1\in(0,1)$, $\rho\in(0,1)$; tolerance $\epsilon>0$
    \Ensure posterior mode $\widehat{\bm{\lambda}}_j$ and curvature $-\mathbf{H}(\widehat{\bm{\lambda}}_j)$
    \State initialize $\bm{\lambda}^{(0)}\gets \mathbf{0}$;\quad $t\gets 0$
    \Repeat
    \State re-profile the offset $\bm{\xi}^{(t)}\gets \bm{\xi}(\bm{\lambda}^{(t)})$
    \State $\mathbf{g}^{(t)}\gets \nabla\ell(\bm{\lambda}^{(t)})$
    \State form $-\mathbf{H}_P(\bm{\lambda}^{(t)})
         =\bar{\mathbf{D}}^{(t)}-\sigma^{-4}\mathbf{B}\mathbf{M}_K^{-1}\mathbf{B}^\top$
    \State solve $\bigl(-\mathbf{H}_P(\bm{\lambda}^{(t)})\bigr)\,\Delta\bm{\lambda}^{(t)}=\mathbf{g}^{(t)}$
    \State $s^{(t)}\gets 1$
    \While{$\ell\bigl(\bm{\lambda}^{(t)}+s^{(t)}\Delta\bm{\lambda}^{(t)}\bigr)
          < \ell(\bm{\lambda}^{(t)})+c_1\,s^{(t)}\,\mathbf{g}^{(t)\top}\Delta\bm{\lambda}^{(t)}$}
    \State $s^{(t)}\gets \rho\,s^{(t)}$
    \EndWhile
    \State $\bm{\lambda}^{(t+1)}\gets \bm{\lambda}^{(t)}+s^{(t)}\Delta\bm{\lambda}^{(t)}$;\quad
         $t\gets t+1$
    \Until{$\lVert\mathbf{g}^{(t)}\rVert\le \epsilon$}
    \State $\widehat{\bm{\lambda}}_j\gets \bm{\lambda}^{(t)}$
    \State \Return $\widehat{\bm{\lambda}}_j$ and
       $-\mathbf{H}(\widehat{\bm{\lambda}}_j)$
    \end{algorithmic}
    \end{algorithm}

\begin{lemma}[Ascent direction]
    Let $\mathbf{g}=\nabla \ell(\bm{\lambda})$ and $\Delta \bm{\lambda}=\bigl(-\mathbf{H}_{P}(\bm{\lambda})\bigr)^{-1}\mathbf{g}$. If $\mathbf{g}\neq \bm{0}$, then $\mathbf{g}^{\top}\Delta \bm{\lambda}>0$.
\end{lemma}

\begin{proof}
    Since $\bm{\Gamma}\prec 0$ implies the Woodbury-expanded matrix in~\ref{eq:inverse-hessian4} is positive definite. Equivalently, \ref{eq:acceleration-hessian3} is positive definite because it is the negative Hessian of the strictly concave function $\ell_{P}(\cdot, \bm{\xi})$. Hence, $-\mathbf{H}_{P}(\bm{\lambda})=\tilde{\mathbf{D}}-\sigma^{-4}\mathbf{BM}_{K}^{-1}\mathbf{B}^{\top}\succ 0$ by Proposition~\ref{prop:negative-definiteness}, so its inverse is also positive definite and $\mathbf{g}^{\top}\Delta\bm{\lambda}=\mathbf{g}^{\top}\bigl(-\mathbf{H}_{P}(\bm{\lambda})\bigr)^{-1}\mathbf{g}>0$ for any $\mathbf{g}\neq 0$.
\end{proof}

Let $\ell: \mathbb{R}^{Q}\to \mathbb{R}$ denote the true conditional log-posterior of the group effect $\bm{\lambda}$, whose mode is defined by~\ref{eq:maximizer}, and $\ell_{P}(\cdot,\bm{\xi})$ denotes the Poisson surrogate log-posterior with offset $\bm{\xi}$. The Newton-type direction is 

\begin{equation}
    \Delta \bm{\lambda}=\bigl(-\mathbf{H}_{P}(\bm{\lambda})\bigr)^{-1}\mathbf{g}, \quad \mathbf{g}=\nabla\ell(\bm{\lambda}),
    \label{eq:newton-type-direction}
\end{equation}
where $\mathbf{H}_{P}(\bm{\lambda})=\nabla^{2}\ell_{P}(\bm{\lambda},\bm{\xi})$ is the Hessian of the Poisson surrogate by the Woodbury expansion~\ref{eq:acceleration-hessian3}. We propose two regularity facts on which the analysis rests, and both of them are consequences of results already established for the model.

\begin{assumption}[Curvature of true objective]
    Let $\ell\in C^{2}(\mathbb{R}^{Q})$. On initial superlevel set $\mathcal{L}_{0}:=\{\bm{\lambda}\in\mathbb{R}^{Q}:\ell(\bm{\lambda})\geq \ell(\bm{\lambda}^{(0)})\}$, there exist constants $0<\mu_{0}\leq L_{0}<\infty$ with
    \begin{equation}
        \mu_{0}\preceq -\nabla^{2}\ell(\bm{\lambda})\preceq L_{0}\mathbf{I},
    \end{equation}
    for all $\bm{\lambda}\in\mathcal{L}_{0}$.
    \label{as:ell}
\end{assumption}

Assumption~\ref{as:ell} holds for the model. The strong-concavity bound is supplied by
the Gaussian random-effect prior alone. Since $\bm{\alpha}\sim N(\mathbf 0,\bm{\Sigma})$ with
$\bm{\Sigma}=\mathbf{B}\mathbf{B}^{\top}+\sigma^{2}\mathbf{I}$, the prior contributes
$\bm{\Sigma}^{-1}\succeq(\theta_{1}+\sigma^{2})^{-1}\mathbf{I}$ to $-\nabla^{2}\ell$, so
$\mu_{0}=(\theta_{1}+\sigma^{2})^{-1}>0$ holds unconditionally.
The smoothness bound is supplied by the multinomial log-likelihood, which has
negative semi-definite curvature and whose operator norm is dominated by the total observed count
$\bigl\lVert\nabla^{2}\log p(\bm{y}\mid\bm{\lambda})\bigr\rVert\le\sum_{i}y_{ij\cdot}<\infty$, so
$L_{0}=\lambda_{\max}(\bm{\Sigma}^{-1})+\sum_{i}y_{ij\cdot}$ is finite whenever the data are finite.
Thus $\ell$ is $\mu_{0}$-strongly concave and $L_{0}$-smooth by construction, and no additional
sampling assumption is needed; the same fact guarantees that the maximizer $\hat{\bm{\lambda}}$
exists and unique, and $\mathcal{L}_{0}$ is compact.

\begin{assumption}[Uniform conditioning of the preconditioner]
    There exist constants $0<m\leq M<\infty$ independent of $t$ such that every preconditioner generated by Algorithm~\ref{alg:estep-inner} satisfies,
    \begin{equation}
        m\mathbf{I}\preceq -\mathbf{H}_{P}(\bm{\lambda}^{(t)})\preceq M\mathbf{I},
    \end{equation}
    for all iteration t.
    \label{as:prec}
\end{assumption}
Assumption~\ref{as:prec} likewise follows from the model, subject to one mild and checkable
condition on the data. The lower bound $m$ is structural: by Proposition~4.2.1 the
Woodbury-expanded matrix~\eqref{eq:woodbury} is positive definite because $\bm{\Gamma}\prec\mathbf 0$,
uniformly over the parameter region of interest. The upper bound $M$ requires that the diagonal
working weights and the profiled offset stay bounded along the iterates; the working weights
already satisfy $d_{q}\le\sum_{i}M_{i}$, so the only quantity that could degrade the conditioning
is the profiled offset $\xi_{ij}$, which enters through $\log y_{ij\cdot}$ and diverges only if a
unit carries no counts at all. It therefore suffices that every included unit has at least one
nonzero count,
\[
y_{ij\cdot}=\sum_{q=1}^{Q}y_{ijq}\ \ge\ 1\qquad\text{for all included }(i,j),
\]
under which $\xi_{ij}$ is finite and, being a continuous function of $\bm{\lambda}$
(Proposition~4.1.1) on the compact super-level set $\mathcal{L}_{0}$, keeps the spectrum of
$-\mathbf{H}_{P}(\bm{\lambda}^{(t)})$ within a fixed interval $[m,M]$. This condition is innocuous
in practice: empty units carry no information and are routinely removed during pre-processing of
count and compositional data, and the same effect is achieved by the standard pseudo-count or
truncation adjustments used for sparse counts. Assumption~\ref{as:prec} thus holds for any dataset
after the elementary pre-processing already standard in this setting.
    
\begin{lemma}[Ascent Direction]
    Let $\mathbf{g}=\nabla\ell(\bm{\lambda})$ and $\Delta \bm{\lambda}=\bigl(-\mathbf{H}_{P}(\bm{\lambda})\bigr)^{-1}\mathbf{g}$ as in~\ref{eq:newton-type-direction}. If $\mathbf{g}\neq \bm{0}$, $\mathbf{g}^{\top}\Delta\bm{\lambda}>0$. Moreover, the angle between $\Delta\bm{\lambda}$ and $\mathbf{g}$ is bounded away from a right angle,
    \begin{equation}
        \frac{\mathbf{g}^{\top}\Delta\bm{\lambda}}{\|\mathbf{g}\|\|\Delta\bm{\lambda}\|}\geq \frac{m}{M}>0.
    \end{equation}
    \label{lem:ascent-direction}
\end{lemma}

\begin{proof}
    By Assumption~\ref{as:prec}, $-\mathbf{H}_{P}(\bm{\lambda})$ is symmetric positive definite, so its inverse is symmetric positive definite as well with spectrum contained in $[M^{-1}, m^{-1}]$. Hence, for any $\mathbf{g}\neq \bm{0}$,
    \begin{equation}
        \mathbf{g}^{\top}\Delta\bm{\lambda}=\mathbf{g}^{\top}\bigl(-\mathbf{H}_{P}(\bm{\lambda})\bigr)^{-1}\mathbf{g}\geq M^{-1}\|\mathbf{g}\|^{2}>0,
    \end{equation}
    which is the first claim. For the angle, using $\mathbf{g}^{\top}\bigl(-\mathbf{H}_{P}(\bm{\lambda})\bigr)^{-1}\mathbf{g}\geq M^{-1}\|\mathbf{g}\|^{2}$ and $\|\Delta\bm{\lambda}\|=\|\bigl(-\mathbf{H}_{P}(\bm{\lambda})\bigr)^{-1}\mathbf{g}\|\leq m^{-1}\|\mathbf{g}\|$,
    \begin{equation}
        \frac{\mathbf{g}^{\top}\Delta\bm{\lambda}}{\|\mathbf{g}\|\|\Delta\bm{\lambda}\|}\geq\frac{M^{-1}\|\mathbf{g}\|^{2}}{m^{-1}\|\mathbf{g}\|^{2}}= \frac{m}{M}>0.
    \end{equation}
\end{proof}

\begin{theorem}[Global Linear Convergence under Armijo Backtracking]
    Let $\bm{\lambda}^{(0)}\in \mathbb{R}^{Q}$ be arbitrary and generate $\bm{\lambda}^{(t+1)}=\bm{\lambda}^{(t)}+s^{(t)}\Delta \bm{\lambda}^{(t)}$, where $\Delta \bm{\lambda}^{(t)}$ is computed from~\ref{eq:newton-direction5} at $\bm{\lambda}^{(t)}$. For fixed Armijo parameters $c_{1}\in(0,1)$ and $\rho\in (0,1)$, the step size $s^{(t)}=\rho^{k_{t}}$ is chosen with $k_{t}\in \{0,1,2,\cdots\}$, which is the smallest integer satisfying the sufficient-increase condition,
    \begin{equation}
        \ell(\bm{\lambda}^{(t)}+s^{(t)}\Delta \bm{\lambda}^{(t)})\geq \ell(\bm{\lambda}^{(t)})+c_{1}s^{(t)}\mathbf{g}^{(t)\top}\Delta \bm{\lambda}^{(t)},
        \label{eq:thm1-1}
    \end{equation}
    where $\mathbf{g}^{(t)}:=\nabla\ell(\bm{\lambda}^{(t)})$. Then the iterates converge to the unique maximizer $\hat{\bm{\lambda}}$ of $\ell$ at a linear rate under Assumption~\ref{as:ell}-\ref{as:prec}. That is, there exists a constant $\kappa\in (0,1]$ depending only on $(\mu_{0},L_{0},m,M,c_{1},\rho)$ such that,
    \begin{equation}
        \ell(\hat{\bm{\lambda}})-\ell(\bm{\lambda}^{(t)})\leq (1-\kappa)^{t}[\ell(\hat{\bm{\lambda}})-\ell(\bm{\lambda}^{(0)})],
        \label{eq:thm1-2}
    \end{equation}
    and,
    \begin{equation}
        \|\hat{\bm{\lambda}}-\bm{\lambda}^{(t)}\|\leq (1-\kappa)^{\frac{t}{2}}\sqrt{\frac{2}{\mu_{0}}[\ell(\hat{\bm{\lambda}})-\ell(\bm{\lambda}^{(0)})]}.
        \label{eq:thm1-3}
    \end{equation}
    \label{thm:convergence}
\end{theorem}

\begin{proof}
    If $\mathbf{g}^{(t)}=\bm{0}$ for some $t$, then $\bm{\lambda}^{(t)}$ is a stationary point of the strictly concave $\ell$, hence $\bm{\lambda}^{(t)}=\hat{\bm{\lambda}}$ and the assertion holds trivially. Under generation assumption, we assume $\mathbf{g}^{(t)}\neq \bm{0}$ and proceed the proof in four steps.
    \medskip

    \noindent\textbf{Step 1.} By Lemma~\ref{lem:ascent-direction}, $\Delta\bm{\lambda}^{(t)}$ is an ascent direction and $\mathbf{g}^{(t)\top}\Delta\bm{\lambda}^{(t)}>0$. Assumption~\ref{as:ell} gives $L_{0}$-Lipschitz continuity of $\nabla\ell$ on $\mathcal{L}_{0}$. Also, the property of monotonicity (proved below) of $\{\ell(\bm{\lambda}^{(t)})\}$ keeps every iterate in $\mathcal{L}_{0}$. Thus, the quadratic lower bound,
    \begin{equation}
        \ell(\bm{\lambda}^{(t)}+s\Delta \bm{\lambda}^{(t)})\geq \ell(\bm{\lambda}^{(t)})+s\mathbf{g}^{(t)\top}\Delta \bm{\lambda}^{(t)}-\frac{L_{0}s^{2}}{2}\|\Delta \bm{\lambda}^{(t)}\|^{2},
        \label{eq:step1-1}
    \end{equation}
    holds for all admissible $s$. Comparing~\ref{eq:step1-1} with~\ref{eq:thm1-1}, the Armijo condition is implied whenever 
    \begin{equation}
        s\leq \frac{2(1-c_{1})\mathbf{g}^{(t)\top}\Delta \bm{\lambda}^{(t)}}{L_{0}\|\Delta \bm{\lambda}^{(t)}\|^{2}}.
        \label{eq:step1-2}
    \end{equation}
    Since backtracking multiplies the trial step by $\rho$ until acceptance and starts from 1, the accepted step satisfies,
    \begin{equation}
        s^{(t)}\geq \min \left\{1,\frac{2\rho(1-c_{1})\mathbf{g}^{(t)\top}\Delta \bm{\lambda}^{(t)}}{L_{0}\|\Delta \bm{\lambda}^{(t)}\|^{2}}\right\}.
        \label{eq:step1-3}
    \end{equation}
    For the simplicity of the notation, we write $P^{(t)}=\bigl(-\mathbf{H}_{P}(\bm{\lambda})\bigr)^{-1}$ with spectrum in $[M^{-1},m^{-1}]$. From Lemma~\ref{lem:ascent-direction}, we have,
    \begin{equation}
        \frac{\mathbf{g}^{(t)\top}\Delta\bm{\lambda}^{(t)}}{\|\Delta\bm{\lambda}^{(t)}\|^{2}}\geq\frac{m^{2}}{M},
        \label{eq:step1-4}
    \end{equation}
    and hence,
    \begin{equation}
        s^{(t)}\geq s_{\min} :=\min \left\{1,\frac{2\rho(1-c_{1})m^{2}}{L_{0}M}\right\}>0,
        \label{eq:step1-5}
    \end{equation}
    which is the bound independent of $t$. In particular, the line search terminates after finitely many backtracking steps.
    \medskip

    \noindent\textbf{Step 2.} Combining the accepted Armijo inequality~\ref{eq:thm1-1} with~\ref{eq:step1-5} and $\mathbf{g}^{\top}\Delta\bm{\lambda}\geq M^{-1}\|\mathbf{g}\|^{2}$, we have,
    \begin{equation}
        \ell(\bm{\lambda}^{(t+1)})-\ell(\bm{\lambda}^{(t)})\geq c_{1}s^{(t)}\mathbf{g}^{(t)\top}\Delta \bm{\lambda}^{(t)} \geq \frac{c_{1}s_{\min}}{M}\|\mathbf{g}^{(t)}\|^{2}.
        \label{eq:step2-1}
    \end{equation}
    We define the positive constant term $\gamma:=\dfrac{c_{1}s_{\min}}{M}$. Thus, $\{\ell(\bm{\lambda}^{(t)})\}$ is non-decreasing and bounded above by $\ell{(\hat{\bm{\lambda})}}$. It converges and every iterate lies in $\mathcal{L}_{0}$, which justify the use of Assumptions~\ref{as:ell}-\ref{as:prec} and of~\ref{eq:thm1-1} above.
    \medskip

    \noindent\textbf{Step 3.} Assumption~\ref{as:ell} yields the strong concavity. That is, for every $\bm{\lambda}\in\mathcal{L}_{0}$, the quadratic upper bound satisfies that
    \begin{equation}
        \ell(\bm{y})\leq \ell(\bm{\lambda})^{\top}(\bm{y}-\bm{\lambda})-\frac{\mu_{0}}{2}\|\bm{y}-\bm{\lambda}\|^{2},
        \label{eq:step3-1}
    \end{equation}
    which maximize the right-hand side over $\bm{y}\in\mathbb{R}^{Q}$. By Polyak-Lojasiewicz inequality, we have,
    \begin{equation}
        \ell(\hat{\bm{\lambda}})-\ell(\bm{\lambda})\leq \frac{1}{2\mu_{0}}\|\nabla\ell(\bm{\lambda})\|^{2}.
        \label{eq:step3-2}
    \end{equation}
    Let $\delta_{t}:=\ell(\hat{\bm{\lambda}})-\ell(\bm{\lambda})\geq 0$. Substracting~\ref{eq:step2-1} from \ref{eq:step3-1} from $\ell(\bm{\lambda})$ and applying \ref{eq:step3-2} at $\bm{\lambda}^{(t)}$,
    \begin{equation}
        \delta_{t+1}=\delta_{t}-\left[\ell(\bm{\lambda}^{(t+1)})-\ell(\bm{\lambda}^{(t)})\right]\leq \delta_{t}-\gamma \|\mathbf{g}^{(t)}\|^{2}\leq \delta_{t}-2\gamma \mu_{0} \delta_{t}=(1-\kappa)\delta_{t},
        \label{eq:step3-3}
    \end{equation}
    where $\kappa:=2\gamma\mu_{0}=\dfrac{2c_{1}s_{\min}\mu_{0}}{M}$. Since $\delta_{t}\geq0$, $\kappa\in(0,1]$ and $\delta_{t}\leq (1-\kappa)^{t}\delta_{0}$, which is the first stated bound and shows $\ell(\hat{\bm{\lambda}})\to\ell(\bm{\lambda})$.
    \medskip

    \noindent\textbf{Step 4.} Strong concavity also gives $\ell(\hat{\bm{\lambda}})\leq\ell(\bm{\lambda})-\dfrac{\mu_{0}}{2}\|\hat{\bm{\lambda}}-\bm{\lambda}\|^{2}$. Hence, 
    \begin{equation}
        \dfrac{\mu_{0}}{2}\|\hat{\bm{\lambda}}-\bm{\lambda}^{(t)}\|^{2}\leq \delta_{t}\leq (1-\kappa)^{t}\delta_{0},
        \label{eq:step4-1}
    \end{equation}
    then,
    \begin{equation}
        \|\hat{\bm{\lambda}}-\bm{\lambda}^{(t)}\|\leq (1-\kappa)^{\frac{t}{2}}\sqrt{\frac{2\delta_{0}}{\mu_{0}}}\to 0.
        \label{eq:step4-2}
    \end{equation}
    Therefore, $\bm{\lambda}^{(t)}\to \hat{\bm{\lambda}}$, which is the unique maximizer of $\ell$.
\end{proof}

\begin{remark}
    The backtracking condition in Theorem~\ref{thm:convergence} is on the true posterior $\ell$ instead of checking on the surrogate $\ell_{P}$. It guarantees that no majorization error from the Poisson approximation accumulates in the final iterate $\hat{\bm{\lambda}}$, which is therefore exactly the mode defined by~\ref{eq:maximizer}, to numerical precision, regardless of how the search direction that got there was computed.
\end{remark}

\section{Curvature Problem \& Hybrid Algorithm}
\label{sec:curvature-problem-hybrid-algorithm}

The posterior mode $\hat{\bm{\lambda}}_{j}$ of the true multinomial log-posterior $\ell$ in E-step for each group $j$. Theorem~\ref{thm:convergence} ganrantees that the Woodbury-accelerated iteration reaches it. However, the Laplace approximation needs the curvature at that mode, and here the two Hessians that the algorithms manipulates part ways. We make the distinction precise, because it is both the source of Woodbury of the speed-up and the source of a small, quantifiable bias.

Recall the linear predictor~\ref{eq:linear-predictor} and the probability $\bm{p}_{ij}=\text{softmax}(\bm{\eta}_{ij})\in\mathbb{R}^{Q}$ for categories. The group effect enters each category additively, where $\dfrac{\partial\eta_{ijq}}{\partial \lambda_{jr}}=\delta_{qr}$, so the score of the true multinomial log-posterior is given by,
\begin{equation}
    \nabla\ell(\bm{\lambda}_{j})=\sum_{i\in\mathcal{I}_{j}}(\bm{y}_{ij}-y_{ij\cdot}\bm{p}_{ij})-\bm{\Sigma}^{-1}\bm{\lambda}_{j},
    \label{eq:true-multi-log-post}
\end{equation}
and the Hessian matrix is given by,
\begin{equation}
    -\mathbf{H}_{j}(\bm{\lambda}_{j})=\sum_{i\in\mathcal{I}_{j}}y_{ij\cdot}\left(\text{diag}(\bm{p}_{ij})-\bm{p}_{ij}\bm{p}_{ij}^{\top}\right)+\bm{\Sigma}^{-1}.
    \label{eq:multi-hessian}
\end{equation}
The prior precision $\bm{\Sigma}^{-1}=\left(\mathbf{BB}^{\top}+\sigma^{2}\mathbf{I}_{Q}\right)^{-1}=\sigma^{-2}\mathbf{I}_{Q}-\sigma^{-4}\mathbf{BM}_{K}^{-1}\mathbf{B}^{\top}$ is diagonal-plus-low-rank-K by the Woodbury identity with $\mathbf{M}_{K}$ defined in~\ref{eq:woodbury2}. However, the likelihood of~\ref{eq:multi-hessian} carries the $I_{j}$ outer products $-y_{ij\cdot}\bm{p}_{ij}\bm{p}_{ij}^{\top}$. It couples all categories together and makes $-\mathbf{H}_{j}$ be a dense $Q\times Q$ matrix. The computational cost $\mathcal{O}(Q^{3})$ prevents the direct Newton solve.

That’s why we introduce the Poisson-Multinomial equivalence at the beginning in Chapter~\ref{cha:model-specification}. The Poisson-Multinomial equivalence represents each count as an independent Poisson variate $y_{ijq}\sim \operatorname{Poisson}\left(\exp (\eta_{ijq}+\xi_{ij})\right)$ with offset $\xi_{ij}=\log y_{ij\cdot}-\log \sum_{q}\exp(\eta_{ijq})$ re-profiled at the current $\bm{\lambda}_{j}$ so that $\exp(\eta_{ijq}+\xi_{ij})=y_{ij\cdot}p_{ijq}$. The surrogate objective $\ell_{P}$ is the Poisson log-posterior with $\xi_{ij}$ held fixed at this profiled value. The score of the surrogate Poisson is,
\begin{equation}
    \nabla\ell_{P}(\bm{\lambda}_{j})=\sum_{i\in\mathcal{I}_{j}}\bigl(\bm{y}_{ij}-\exp (\eta_{ijq}+\xi_{ij})\bigr)-\bm{\Sigma}^{-1}\bm{\lambda}_{j}=\sum_{i\in\mathcal{I}_{j}}(\bm{y}_{ij}-y_{ij\cdot}\bm{p}_{ij})-\bm{\Sigma}^{-1}\bm{\lambda}_{j},
    \label{eq:poi-log-post}
\end{equation}
coincides with the true multinomial score in~\ref{eq:true-multi-log-post}. The surrogate and the target share the same stationary point, which is why the line search recovers the exact multinomial model. The Hessians differ. Holding $\xi_{ij}$ fixed makes each Poisson coordinate a function of its own $\lambda_{jq}$ alone, so the cross-category second derivatives vanish and,
\begin{equation}
    -\mathbf{H}_{P}(\bm{\lambda}_{j})=\text{diag}\left(\sum_{i\in\mathcal{I}_{j}}y_{ij\cdot}\bm{p}_{ij}\right)+\bm{\Sigma}^{-1}=\tilde{\mathbf{D}}-\sigma^{-4}\mathbf{BM}_{K}^{-1}\mathbf{B}^{\top},
    \label{eq:poi-hessian}
\end{equation}
where $\tilde{d}_{q}=\sum_{i\in\mathcal{I}_{j}}y_{ij\cdot}p_{ijq}+\sigma^{-2}$. Therefore, the surrogate negative Hessian is diagonal-plius-rank-K, which is the structure solved by Woodbury equality and gives the Newton direction $\left(-\mathbf{H}_{P}\right)^{-1}\nabla\ell$ and the log-determinant $\log |-\mathbf{H}_{P}|$ with cost $\mathcal{O}(QK^{2}+K^{3})$ instead of $\mathcal{O}(Q^{3})$.

\begin{proposition}
    For every $\bm{\lambda}_{j}$,
    \begin{equation}
        -\mathbf{H}_{j}(\bm{\lambda})=-\mathbf{H}_{P}(\bm{\lambda})-\sum_{i\in \mathcal{I}_{j}}y_{ij\cdot} \bm{p}_{ij}\bm{p}_{ij}^{\top}.
        \label{eq:diff-hessian}
    \end{equation}
    The correction is positive semi-definite of rank at most $I_{j}=|\mathcal{I}_{j}|$, consequently $-\mathbf{H}_{P}(\bm{\lambda})\succeq -\mathbf{H}_{j}(\bm{\lambda})\succ 0$.
\end{proposition}

\begin{proof}
    Subtract \ref{eq:poi-hessian} from \ref{eq:multi-hessian}, the diagonal terms $\text{diag}\left(\sum_{i}y_{ij\cdot}\bm{p}_{ij}\right)$ and the prior precision $\bm{\Sigma}^{-1}$ cancel, leaving the equality \ref{eq:diff-hessian}. Each summand $y_{ij\cdot}\bm{p}_{ij}\bm{p}_{ij}^{\top}$ is a non-negative scalar times an outer product, hence positive semi-definite of rank one. Their sum has rank at most $I_{j}$. Positive semi-definiteness of the correction gives $-\mathbf{H}_{P}\succeq -\mathbf{H}_{j}$, and $-\mathbf{H}_{j}\succ 0$ because it is the negative Hessian of the strictly concave $\ell$. 
\end{proof}

so that 
\begin{equation}
    \hat{\mathbf{S}}_{j}^{P}:=\bigl(-\mathbf{H}_{P}(\hat{\bm{\lambda}_{j}})\bigr)_{-1}\preceq \hat{\mathbf{S}}_{j}, 
\end{equation}
in the positive semi-definite order. The Poisson surrogate systematically understates posterior uncertainty. Substituting $\hat{\mathbf{S}}_{j}^{P}$ for $\hat{\mathbf{S}}_{j}$ in~\ref{eq:sigma-obs} would therefore introduce a systematic downward bias into $\hat{\sigma}^{2}$ and into $\mathcal{L}^{(t)}$.

The resolution is the hybrid strategy that gives the name of this section. The cheap Woodbury-accelerated surrogate Newton iteration is used to find the model $\hat{\bm{\lambda}}_{j}$ and evaluate the true multinomial Hessian once to obtain $\hat{\mathbf{S}}_{j}$ at that mode.

Theorem~\ref{thm:convergence} guarantees that the Woodbury-accelerated iteration converges to the correct mode $\hat{\bm{\lambda}}$. It says noting, however, about the curvature at that mode, and the curvature is what the Laplace approximation~\ref{eq:surrogate-eq} actually needs

\begin{equation}
    \hat{\mathbf{S}}_{j}=\bigl(-\mathbf{H}_{j}(\hat{\bm{\lambda}}_{j})\bigr)^{-1},
\end{equation}
with $\mathbf{H}_{j}$ the Hessian of the true multinomial log-posterior $\ell$, not of the Poisson surrogate $\ell_{P}$. 

\iffalse
\section{Complexity Analysis}
\label{sec:complexity-analysis}

\section{Woodbury-Accelerated Factor-Analysis M-step}
\label{sec:woodbury-accelerated-factor-analysis}
\fi

\section{Fixed-Q Consistency}
\label{sec:fixed-q-consistency}

Before turning to simulation evidence, we record a consistency result in the classical asymptotic regime $J\to \infty$ with $Q$ and $K$ fixed, which confirms that the Laplace-EM procedure of Chapter~\ref{cha:methodology} is well-defined. We give the model assumptions for the proof. 

\begin{assumption}[Identifiability]
    The PLT constraint~\ref{eq: PLT} holds,
\end{assumption}

\begin{assumption}[Bounded Moments]
    $\mathbb{E}[y_{ij\cdot}]<\infty$ for any observation $i$,
\end{assumption}

\begin{assumption}[Compact Parameter Space]
    Parameter space $\bm{\theta}=\{\bm{\mu},\mathbf{B},\sigma^{2},\bm{\Phi}\}\in \mathcal{F}$ with $\mathcal{F}$ compact and $\sigma^{2}\geq \sigma^{2}_{\text{min}}>0$,
\end{assumption}

\begin{assumption}[Non-degenerate design]
    $\inf_{j}\lambda_{\text{min}}(\mathbf{X}_{j}^{\top}\mathbf{X}_{j})>c>0$, where $\mathbf{X}_{j}$ collects the covariate vectors of group $j$.
\end{assumption}

\begin{theorem}[Fixed-Q consistency]
    Under assumption and standard smoothness conditions on the log-likelihood, as $J\to\infty$ with $Q$, $K$ and within-group sample size $N_{j}$ fixed, the EM stationary point $\hat{\bm{\theta}}_{J}$ satisfies that $\hat{\bm{\theta}}_{J}\to \bm{\theta}^{*}$, $\bm{\theta}^{*}$ denotes the true value.
    \label{thm:fixed-q-theorem}
\end{theorem}

\begin{proof}
    
\end{proof}

Theorem~\ref{thm:fixed-q-theorem} is reassuring but limited. In this theorem, we fix $Q$ and let $J\to\infty$, whereas the scientific motivation of this thesis (Chapter~\ref{cha:intro}) is precisely the opposite regime. In the next chapter we will discuss the Neyman-Scott problems in our case and the relation of estimator and true value.
